\documentclass[11pt]{article}

\usepackage[margin=1in]{geometry}
\usepackage{amsmath, amssymb}
\usepackage{xcolor}
\usepackage[skins,breakable]{tcolorbox}
\usepackage{hyperref}

\hypersetup{
    colorlinks=true,
    linkcolor=blue!60!black,
    urlcolor=blue!60!black,
    citecolor=blue!60!black
}

\newcommand{\SymPy}{SymPy}

% Public artifact repository; \bugbundle links a bug's bundle folder into it.
% TODO: set the public artifact repository URL before release.
\newcommand{\artifactrepo}{https://github.com/TODO-org/TODO-repo}
\newcommand{\bugbundle}[1]{\href{\artifactrepo/tree/main/bug-report/#1}{\nolinkurl{bug-report/#1/}}}

\newtcolorbox{counterexamplebox}{
  enhanced,
  colback=yellow!5,
  colframe=orange!70!black,
  arc=2mm,
  boxrule=0.8pt,
  left=4mm,
  right=4mm,
  top=3mm,
  bottom=3mm,
  before=\par\medskip\noindent,
  after=\par\medskip,
  before upper=\raggedright
}

\title{SymPy Bug Card}
\date{}

\begin{document}

\author{}
\maketitle

\section*{Bug 52: \texttt{function\_range} Says \texttt{sqrt(tan(x)**2)} Is Constant}

\begin{counterexamplebox}
\textbf{Task.} Compute the real range of the expression on the interval
\([-1,1]\).
\[
  \operatorname{function\_range}\!\left(\sqrt{\tan(x)^2}, x, [-1,1]\right)
\]

\medskip
\textbf{SymPy's answer.}
\[
  \operatorname{function\_range}\!\left(\sqrt{\tan(x)^2}, x, [-1,1]\right)
  = \{\tan(1)\}.
\]

\medskip
\textbf{Correct answer.}
\[
  [0,\tan(1)].
\]

\medskip
\textbf{Explanation.} On \([-1,1]\), \(\sqrt{\tan(x)^2}=|\tan(x)|\). The
function is continuous there, has value \(0\) at \(x=0\), and has maximum
\(\tan(1)\) at both endpoints, so the singleton answer omits attained values.

\medskip
\textbf{Likely root cause.} The diagnosis points to
\nolinkurl{sympy/calculus/util.py}: \texttt{function\_range} samples endpoints
and derivative zeros, but misses points where the derivative is undefined while
the original function remains continuous. This is a new instance of a known
failure mode.

\medskip
\textbf{Artifact (GitHub).} \bugbundle{bug-052-function-range-says-sqrt-tan-x-2-is-constant}
\end{counterexamplebox}

\end{document}
